
% ---------- 字体与数学排版 ----------
\usepackage{fontspec}
\usepackage{xeCJK}
\usepackage{mathtools}
\usepackage{unicode-math}
\setmainfont{STIXTwoText-Regular.otf}[
  BoldFont=STIXTwoText-Bold.otf,
  ItalicFont=STIXTwoText-Italic.otf,
  BoldItalicFont=STIXTwoText-BoldItalic.otf
]
\setsansfont{STIXTwoText-Regular.otf}[
  BoldFont=STIXTwoText-Bold.otf,
  ItalicFont=STIXTwoText-Italic.otf,
  BoldItalicFont=STIXTwoText-BoldItalic.otf
]
\setmonofont{NotoSansMono-Regular.ttf}[
  BoldFont=NotoSansMono-Bold.ttf,
  ItalicFont=NotoSansMono-Regular.ttf,
  BoldItalicFont=NotoSansMono-Bold.ttf,
  ItalicFeatures={FakeSlant=0.2},
  BoldItalicFeatures={FakeSlant=0.2}
]
\setmathfont{STIXTwoMath-Regular.otf}
\setCJKmainfont{FandolFang-Regular.otf}[
  AutoFakeBold=2.1,
  AutoFakeSlant=0.15
]
\setCJKsansfont{FandolFang-Regular.otf}[
  AutoFakeBold=2.1,
  AutoFakeSlant=0.15
]
\setCJKmonofont{FandolFang-Regular.otf}[
  AutoFakeBold=2.0,
  AutoFakeSlant=0.15
]

% ---------- 页面、图表与结构 ----------
\usepackage[letterpaper,top=2.25cm,bottom=2.25cm,left=2.35cm,right=2.35cm,headheight=28pt]{geometry}
\usepackage{graphicx}
\usepackage{booktabs}
\usepackage{tabularx}
\usepackage{longtable}
\usepackage{array}
\usepackage{multirow}
\usepackage{makecell}
\usepackage{enumitem}
\usepackage{siunitx}
\usepackage[table]{xcolor}
\usepackage{tikz}
\usetikzlibrary{arrows.meta,positioning,calc,fit,matrix,backgrounds}
\usepackage[most]{tcolorbox}
\usepackage{listings}
\usepackage{caption}
\usepackage{subcaption}
\usepackage{float}
\usepackage{fancyhdr}
\usepackage{titlesec}
\usepackage{hyperref}
\usepackage{bookmark}
\usepackage{url}
\usepackage{lastpage}
\usepackage{setspace}
\usepackage{needspace}
\usepackage{placeins}
\usepackage{etoolbox}
\usepackage{import}
\raggedbottom
\pretocmd{\section}{\Needspace{10\baselineskip}}{}{}
\pretocmd{\subsection}{\Needspace{10\baselineskip}}{}{}
\pretocmd{\subsubsection}{\Needspace{5\baselineskip}}{}{}

% ---------- 颜色 ----------
\definecolor{UIUCBlue}{HTML}{13294B}
\definecolor{UIUCOrange}{HTML}{FF5F05}
\definecolor{Ink}{HTML}{17202A}
\definecolor{SoftBlue}{HTML}{EEF3F8}
\definecolor{SoftOrange}{HTML}{FFF1E8}
\definecolor{SoftGray}{HTML}{F4F5F6}
\definecolor{DeepGreen}{HTML}{1B6B57}
\definecolor{SoftGreen}{HTML}{EAF5F1}
\definecolor{WarnRed}{HTML}{9B2C2C}
\definecolor{SoftRed}{HTML}{FFF0F0}
\definecolor{CodeBg}{HTML}{F7F7F8}
\definecolor{CodeFrame}{HTML}{CBD2D9}
\definecolor{CodeGray}{HTML}{57606A}
\colorlet{UIUCNavy}{UIUCBlue}

% ---------- 全局排版 ----------
\setstretch{1.34}
\setlength{\parindent}{2em}
\setlength{\parskip}{0.35em}
\setlength{\emergencystretch}{2em}
\setlist[itemize]{leftmargin=2.2em,itemsep=0.3em,topsep=0.35em}
\setlist[enumerate]{leftmargin=2.35em,itemsep=0.35em,topsep=0.35em}
\captionsetup{font=small,labelfont=bf,labelsep=quad,justification=centering}
\sisetup{detect-all,group-separator={,},group-minimum-digits=4}
\urlstyle{same}
\setcounter{tocdepth}{2}
\setcounter{secnumdepth}{3}

% ---------- 标题 ----------
\ctexset{
  section={
    format=\Large\bfseries\color{UIUCBlue},
    beforeskip=1.2em,
    afterskip=0.65em
  },
  subsection={
    format=\large\bfseries\color{Ink},
    beforeskip=1em,
    afterskip=0.45em
  },
  subsubsection={
    format=\normalsize\bfseries\color{black},
    beforeskip=0.8em,
    afterskip=0.35em
  }
}
\titleformat{\section}
  {\Large\bfseries\color{UIUCBlue}}
  {\tikz[baseline=(n.base)]\node[fill=UIUCOrange,text=white,rounded corners=1.5pt,inner xsep=6pt,inner ysep=2pt] (n) {\thesection};}
  {0.65em}{}

% ---------- 页眉页脚与链接 ----------
\pagestyle{fancy}
\fancyhf{}
\fancyhead[L]{\small ECE 408}
\newcommand{\NotesRunningTitle}[1]{\fancyhead[R]{\parbox[b]{0.78\textwidth}{\raggedleft\scriptsize#1}}}
\NotesRunningTitle{\nouppercase{\leftmark}}
\newif\ifNotesFrontmatter
\newcommand{\NotesPageInfo}{第 \thepage\ 页\ifNotesFrontmatter\else，共 \pageref{LastPage}\ 页\fi}
\fancyfoot[L]{\small Applied Parallel Programming}
\fancyfoot[R]{\small\NotesPageInfo}
\renewcommand{\headrulewidth}{0.4pt}
\renewcommand{\footrulewidth}{0.2pt}
\fancypagestyle{plain}{\fancyhf{}\fancyfoot[L]{\small Applied Parallel Programming}\fancyfoot[R]{\small\NotesPageInfo}\renewcommand{\headrulewidth}{0pt}\renewcommand{\footrulewidth}{0.2pt}}
\hypersetup{colorlinks=true,linkcolor=UIUCBlue,urlcolor=DeepGreen,citecolor=UIUCOrange,bookmarksopen=true,bookmarksopenlevel=1,bookmarksdepth=3}

% ---------- 常用命令 ----------
\newcommand{\term}[2]{\textbf{#1}\,(\textit{#2})}
\newcommand{\code}[1]{\texttt{\detokenize{#1}}}
\newcommand{\important}[1]{\textcolor{WarnRed}{\textbf{#1}}}
\colorlet{Navy}{UIUCBlue}
\colorlet{Orange}{UIUCOrange}
\colorlet{Muted}{CodeGray}
\colorlet{Pale}{SoftBlue}
\colorlet{Rule}{CodeFrame}
\colorlet{PaleBlue}{SoftBlue}
\colorlet{RuleGray}{CodeFrame}
\colorlet{UIUCblue}{UIUCBlue}
\colorlet{UIUCorange}{UIUCOrange}
\colorlet{InkGray}{CodeGray}
\newcommand{\Bunit}{\,\mathrm{B}}
\newcommand{\GBs}{\,\mathrm{GB/s}}
\newcommand{\T}{T}

\newcolumntype{Y}{>{\raggedright\arraybackslash}X}
\newcolumntype{C}[1]{>{\centering\arraybackslash}p{#1}}
\newcolumntype{L}[1]{>{\raggedright\arraybackslash}p{#1}}

% ---------- 信息框 ----------
\newtcolorbox{overviewbox}[1][]{
  enhanced,breakable,
  colback=SoftBlue,colframe=UIUCBlue,
  boxrule=0.75pt,arc=2mm,
  left=1.2mm,right=1.2mm,top=1mm,bottom=1mm,
  title={内容概览},fonttitle=\bfseries,
  #1
}
\newtcolorbox{definitionbox}[1][]{
  enhanced,breakable,
  colback=SoftBlue,colframe=UIUCBlue,
  boxrule=0.75pt,arc=2mm,
  left=1.2mm,right=1.2mm,top=1mm,bottom=1mm,
  title={定义},fonttitle=\bfseries,
  #1
}
\newtcolorbox{keybox}[1][]{
  enhanced,breakable,
  colback=SoftGreen,colframe=DeepGreen,
  boxrule=0.7pt,arc=1.5mm,
  left=1.2mm,right=1.2mm,top=1mm,bottom=1mm,
  title={关键结论},fonttitle=\bfseries,
  #1
}
\newtcolorbox{mechanismbox}[1][]{
  enhanced,breakable,
  colback=SoftGreen,colframe=DeepGreen,
  boxrule=0.7pt,arc=1.5mm,
  left=1.2mm,right=1.2mm,top=1mm,bottom=1mm,
  title={机制},fonttitle=\bfseries,
  #1
}
\newtcolorbox{warningbox}[1][]{
  enhanced,breakable,
  colback=SoftRed,colframe=WarnRed,
  boxrule=0.7pt,arc=1.5mm,
  left=1.2mm,right=1.2mm,top=1mm,bottom=1mm,
  title={需要特别注意},fonttitle=\bfseries,
  #1
}
\newtcolorbox{examplebox}[1][]{
  enhanced,breakable,
  colback=SoftOrange,colframe=UIUCOrange,
  boxrule=0.7pt,arc=1.5mm,
  left=1.2mm,right=1.2mm,top=1mm,bottom=1mm,
  title={例题},fonttitle=\bfseries,
  #1
}
\newtcolorbox{derivationbox}[1][]{
  enhanced,breakable,
  colback=SoftGray,colframe=gray!65,
  boxrule=0.65pt,arc=1.5mm,
  left=1.2mm,right=1.2mm,top=1mm,bottom=1mm,
  title={推导补充},fonttitle=\bfseries,
  #1
}

% ---------- CUDA 代码 ----------
\lstdefinestyle{cuda}{
  language=C++,
  morekeywords={__global__,__device__,__host__,__shared__,__syncthreads,dim3,threadIdx,blockIdx,blockDim,gridDim,warpSize,cudaGetDeviceProperties,cudaDeviceProp,cudaError_t},
  basicstyle=\ttfamily\small\linespread{1}\selectfont,
  keywordstyle=\color{UIUCBlue}\bfseries,
  commentstyle=\color{DeepGreen}\itshape,
  stringstyle=\color{WarnRed},
  numbers=left,
  numberstyle=\scriptsize\color{gray},
  numbersep=8pt,
  frame=single,
  rulecolor=\color{CodeFrame},
  backgroundcolor=\color{CodeBg},
  showstringspaces=false,
  breaklines=true,
  columns=fullflexible,
  keepspaces=true,
  tabsize=4,
  xleftmargin=1.2em,
  framexleftmargin=0.7em
}
\lstset{style=cuda}
\renewcommand{\lstlistingname}{代码}
\renewcommand{\lstlistlistingname}{代码目录}


% ---------- 兼容各讲的数学、图片与信息框 ----------
\newcommand{\platform}[2]{#1\,(\textit{#2})}
\newcommand{\vect}[1]{\boldsymbol{#1}}
\newcommand{\ceildiv}[2]{\left\lceil\dfrac{#1}{#2}\right\rceil}
\newcommand{\occupancy}{\mathcal{O}}
\newtcolorbox{modelbox}[1][]{
  enhanced,breakable,colback=SoftGray,colframe=gray!65,
  boxrule=0.65pt,arc=1.5mm,
  left=1.2mm,right=1.2mm,top=1mm,bottom=1mm,
  title={建模说明},fonttitle=\bfseries,#1
}
% graphicspath 由入口设置；缺图直接报错，不生成占位框。
\newcommand{\slidefig}[3]{%
  \begin{figure}[H]\centering
    \includegraphics[width=#1\textwidth]{#2}
    \caption{#3}
  \end{figure}
}

% ---------- 单讲与合订本共用的元数据 ----------
\newcommand{\DeclareLecture}[3]{%
  \ifcsdef{FirstLecture}{}{\gdef\FirstLecture{#1}}%
  \gdef\LastLecture{#1}%
  \expandafter\gdef\csname eceTitle#1\endcsname{#2}%
  \expandafter\gdef\csname eceKeywords#1\endcsname{#3}%
  \listgadd{\AllLectures}{#1}%
}
\newcommand{\LectureTitle}[1]{\csname eceTitle#1\endcsname}
\newcommand{\LectureKeywords}[1]{\csname eceKeywords#1\endcsname}
\newcommand{\LectureRange}{\FirstLecture--\LastLecture}
\newcommand{\LectureCompatibility}[1]{%
  % L1--3 的行内代码使用 LaTeX 转义字符；L4 起使用字面量。
  \ifnum#1<4
    \renewcommand{\code}[1]{\texttt{##1}}%
  \else
    \renewcommand{\code}[1]{\texttt{\detokenize{##1}}}%
  \fi
  \lstset{style=cuda}%
}

% ---------- 共用封面 ----------
\newcommand{\NotesCover}[2]{%
  \hypersetup{pageanchor=false}%
  \begin{titlepage}
    \thispagestyle{empty}
    \begin{tikzpicture}[remember picture,overlay]
      \ifstrempty{#1}{%
        % 合订本以课程编号为唯一标题。
        \fill[UIUCBlue] (current page.north west) rectangle ([yshift=-10cm]current page.north east);
        \node[anchor=west,inner sep=0pt,text=white]
          at ([xshift=2.6cm,yshift=-6.5cm]current page.north west)
          {{\fontsize{56}{64}\selectfont\bfseries ECE 408}};
        \fill[UIUCOrange] ([xshift=1.8cm,yshift=-10cm]current page.north west) rectangle ([xshift=2.15cm,yshift=-12.3cm]current page.north west);
      }{%
        % 单讲保留课程编号、讲次与技术标题。
        \fill[UIUCBlue] (current page.north west) rectangle ([yshift=-4.2cm]current page.north east);
        \node[anchor=west,inner sep=0pt,text=white]
          at ([xshift=2.35cm,yshift=-2.4cm]current page.north west)
          {{\fontsize{28}{34}\selectfont\bfseries ECE 408}};
        \fill[UIUCOrange] ([xshift=1.8cm,yshift=-4.2cm]current page.north west) rectangle ([xshift=2.15cm,yshift=-6.5cm]current page.north west);
      }
      \draw[UIUCOrange,line width=2.2pt] ([xshift=1.8cm,yshift=2.0cm]current page.south west) -- ([xshift=-1.8cm,yshift=2.0cm]current page.south east);
    \end{tikzpicture}
    \ifstrempty{#1}{}{%
      \vspace*{4.8cm}
      \begin{flushleft}
        \hspace*{1.25cm}\begin{minipage}{\dimexpr\linewidth-1.25cm\relax}
          {\Huge\bfseries\color{UIUCBlue}#1\par}
          \vspace{0.5cm}
          {\raggedright\hyphenpenalty=10000\LARGE\bfseries\color{Ink}#2\par}
        \end{minipage}
      \end{flushleft}
    }
    \vfill
  \end{titlepage}
}
\newcommand{\NotesContents}{%
  \NotesFrontmattertrue
  \begingroup
    \small\setstretch{1.12}\setlength{\parskip}{0pt}
    \tableofcontents
  \endgroup
  \clearpage
  \NotesFrontmatterfalse
}
\newcommand{\StandaloneLecture}[1]{%
  \LectureCompatibility{#1}%
  \NotesRunningTitle{Lecture #1: \LectureTitle{#1}}%
  \graphicspath{{assets/}}%
  \hypersetup{pdftitle={ECE408 Lecture #1: \LectureTitle{#1}},pdfsubject={Applied Parallel Programming Lecture #1 Notes},pdfkeywords={\LectureKeywords{#1}}}%
  \NotesCover{第#1讲}{\LectureTitle{#1}}%
  \pagenumbering{roman}\hypersetup{pageanchor=true}%
  \markboth{Lecture #1: \LectureTitle{#1}}{}%
  \NotesContents
  \pagenumbering{arabic}%
  \input{content.tex}%
}
